% T21 sources: limitations.tex; threats_to_validity.tex
The formal model assumes that the directed acyclic graph $G=(V,E)$ is known
before planning begins and remains unchanged during execution. This assumption
is reasonable for a prespecified workload but limits the model's use for
research tasks, architectural explorations, and other R\&D
processes in which new tasks and dependencies are discovered as results become
available. In such processes, the estimates $L_4(P)$ and $B_4$ apply only to the
current snapshot $G_t$; a critical path computed once cannot be interpreted as
a fixed property of the entire project.

All tasks are assumed to be available when planning begins, their phases are
assumed to be nonpreemptive, and the only constrained resources are the $P$
agent streams and one developer. Release dates, separate CI/API capacities, and
fully manual tasks require a different formulation; M4 describes only a
prespecified AI-assisted task set.

The model also does not describe the mechanism by which new vertices are
revealed, the cost of replanning, or the choice of how to decompose the work
further. Splitting a task can expose only substantively feasible parallelism: it
does not eliminate genuine precedence relations and usually increases the cost
of task specification, review, and integration. The reported results can
therefore be used to evaluate a given graph or its current snapshot, but do not
by themselves determine an optimal policy for constructing a dynamic graph.

An absolute prediction of completion time requires calibrating $X$, $k_i$,
$h_i$, $C(P)$, and $\gamma(P)$ for the specific tools and operating mode; when
these change, the calibration may quickly become obsolete. The requirements
are weaker for a comparative choice between two variants: workloads and
durations may be expressed in notional units of a single baseline task if both
configurations are evaluated on a common scale.
Proposition~\ref{prop:invariance} guarantees ranking robustness only to a single
common multiplier applied equally to all agent and human-attention phases of
the variants being compared. Unequal errors across tasks or modes, distortion
of the decomposition between $a_i$ and $h_i$ or of the within-task phase
profile, and different values of $C(P)$ and $\gamma(P)$ can change the resource
paths and the ranking of the variants.

\textbf{Construct validity.}
The coefficients $k$, $a$, and $h$ aggregate the quality of task specification,
the capabilities of the tool, and the review and correction of the result. The
same numerical profile may conceal different mechanisms, and makespan does not
measure code quality, understanding of the system, computational cost, or
downstream defects. The terms ``synchronous'' and ``concentration'' in the early
experimental hypotheses proved insufficiently operationalized; the associated
conclusions are not treated as supported.

\textbf{Internal validity.}
The exact series use deterministic durations, an integer grid, and a single
solver. The risk of implementation error is reduced through independent checks
of the constraints, optimality statuses, frozen inputs, and reproducible
manifests; however, a computational experiment verifies properties of the
implementation of the formal model, not the truth of its assumptions about the
real process.

\textbf{Conclusion validity.}
Finite exact grids do not prove universal patterns beyond the selected ranges.
The counts 30/90, 12/13, and 96/100 describe the matrices examined, not
frequencies in the population of projects. In the large EXP-07 series and part
of EXP-08, a fixed policy is evaluated; its effects cannot be transferred to
$T^*$ without a separate exact proof.

\textbf{External validity.}
Most of the DAGs, durations, and functions $C(P),\gamma(P)$ are synthetic and
have not been calibrated on longitudinal data from teams. The formulation
considers one developer, homogeneous agent streams, and local blocking with
slot retention. Other schedulers may reuse capacity held by a waiting task,
multiple people may handle requests in parallel, and real tasks may reveal
dependencies and rework dynamically.

\textbf{Parameter uncertainty.}
EXP-06 tests only prespecified random and adverse error families. The observed
robustness of the ranking does not guarantee robustness to an arbitrary joint
error in the estimates. Before a practical decision is made, local calibration,
interval analysis, and replanning are required as actual data become available.
